\chapter{Neural Operators}
\label{ch:operators}

\epigraph{\itshape Operator learning is to function approximation what
calculus is to arithmetic.}{Anandkumar, Stuart and collaborators}

\section{Solution operators}
The forward map of a parameterised PDE family is an operator
\begin{equation}
\Gop : a\in\mathcal{A} \;\longmapsto\; u\in\mathcal{U},
\end{equation}
between Banach spaces of functions. Where a PINN learns one $u$, a
\emph{neural operator} learns $\Gop$. Once trained, evaluating $\Gop$
on a new $a$ is a forward pass --- orders of magnitude faster than a
fresh PDE solve.

\section{DeepONet}
DeepONet~\cite{lu2021deeponet} factorises the operator into branch and
trunk networks:
\begin{equation}
\Gop[a](\bm{y}) \approx \sum_{k=1}^p b_k(a(\bm{x}_1),\dots,a(\bm{x}_m))\,t_k(\bm{y}),
\end{equation}
where the \emph{branch} $\bm{b}$ encodes the input function (sampled at
sensors $\bm{x}_i$) and the \emph{trunk} $\bm{t}$ encodes the query
point $\bm{y}$. The architecture is justified by a universal
approximation theorem for nonlinear continuous operators (Chen \& Chen
1995). Variants: POD-DeepONet, multiple-input DeepONet, physics-informed
DeepONet.

\section{Fourier Neural Operator (FNO)}
\label{sec:fno}
FNO~\cite{li2021fno} parameterises the operator as a sequence of layers
\begin{equation}
v_{\ell+1}(\bm{x}) = \sigma\!\left(\bm{W}v_\ell(\bm{x}) + \F^{-1}(\bm{R}_\ell\cdot \F v_\ell)(\bm{x})\right),
\end{equation}
where $\F$ is the Fourier transform and $\bm{R}_\ell$ is a learned
truncated convolution kernel acting on the lowest $K$ modes. Key
properties:
\begin{itemize}
\item Discretisation invariance: trained on one resolution, queried at
another (with caveats).
\item Quasi-linear cost via FFT: $\Oh(N\log N)$ per layer.
\item Strong empirical performance on smooth periodic problems.
\end{itemize}

\subsection{Variants of FNO}
\paragraph{Geo-FNO, Factorised FNO, Tensor-FNO.}
Handle non-rectangular domains by warping geometries; factorise modes
across dimensions; tensor-decompose the kernel.
\paragraph{U-FNO, U-Net FNO.}
Combine FNO with U-Net hierarchy for multiscale features.
\paragraph{Spherical FNO}~\cite{bonev2023sfno}.
Replace FFT with spherical harmonic transforms; used for
state-of-the-art global weather emulators (FourCastNet, GraphCast
hybrids).
\paragraph{Adaptive FNO and AFNO.}
Adapt mode truncation per token, used inside vision transformers.

\begin{pitfall}{FNO discretisation invariance, with asterisks}
Out-of-distribution behaviour at resolutions far from training is
dataset-dependent. Aliasing, periodic-padding artefacts, and limited
mode budgets can all break the ``train once, query anywhere'' promise.
Always benchmark on multiple resolutions.
\end{pitfall}

\section{Graph Neural Operators}
For unstructured meshes, GNOs~\cite{li2020gno,brandstetter2022mpneural}
parameterise an integral operator
\begin{equation}
v_{\ell+1}(\bm{x}) = \sigma\!\left(\bm{W}v_\ell(\bm{x}) + \int_{B_r(\bm{x})}\kappa_\theta(\bm{x},\bm{y})v_\ell(\bm{y})\,\dd\bm{y}\right)
\end{equation}
via message passing on a graph built from the mesh. MeshGraphNet,
GNN-FNO hybrids, and the message-passing PDE solver belong to this
family.

\section{Transformer-based operators}
Transformers naturally handle variable-length, irregular inputs. Recent
high-impact systems include:
\begin{itemize}
\item \textbf{GNOT, OFormer, GAOT}: cross-attention between query
points and conditioning.
\item \textbf{Transolver}: token-mixing in physical space with attention
across slices.
\item \textbf{DPOT, MPP, Poseidon}: transformer-based foundation models
(\cref{ch:emerging}).
\end{itemize}

\paragraph{Poseidon}~\cite{herde2024poseidon} is a multiscale operator
transformer pretrained on a heterogeneous mix of PDE families with
\emph{time-conditioned layer norms} and a semi-group training loss; it
demonstrates strong few-shot transfer to new PDEs.

\section{Discretisation invariance: theory}
A defining requirement of a \emph{neural operator} (vs.\ a
discretisation-tied surrogate) is that it is well-defined as a map
between function spaces, with finite-dimensional implementations
recovered as projections~\cite{kovachki2023neural}. Universal
approximation theorems exist for FNO, DeepONet, and graph operators
under appropriate compactness assumptions.

\section{Training data}
\begin{recipe}{Building a neural-operator dataset}
\begin{enumerate}
\item Define the PDE family by sampling parameters/coefficients/ICs from
a distribution that covers downstream use.
\item Solve with a high-fidelity classical solver; store inputs $a_i$
and solutions $u_i$ at multiple resolutions / time steps.
\item Decide loss: relative $L^2$, $H^1$, or task-specific energy norms.
\item For time-dependent operators, choose autoregressive (next-step) or
direct (input-to-final) framing; combine with rollout regularisation.
\item Augment with symmetries (translation, rotation, reflection) when
the PDE allows.
\end{enumerate}
\end{recipe}

\section{Beyond pure data: physics-informed operators}
Add residual losses on the operator output:
\begin{equation}
\Lop_{\text{op}}=\E_{a}\big[\|\Gop_\theta[a]-u_a\|^2 + \lambda\|\Nop_a[\Gop_\theta[a]]\|^2\big].
\end{equation}
PI-DeepONet, PINO, and FNO-PINN hybrids report improved generalisation
and faster convergence with limited data.

\section{Limitations and open problems}
\begin{itemize}
\item Generalisation to unseen geometries / boundary conditions.
\item High-frequency / shock fronts: most operators are biased to
smooth solutions.
\item Long-horizon temporal stability: errors accumulate in
autoregressive rollout.
\item Foundation-model scaling: data heterogeneity is harder than in
language; tokenisation of physical fields is unsolved.
\end{itemize}

\section{Choosing an operator architecture}
\begin{insight}{Quick decision guide}
\begin{itemize}
\item Periodic / regular grid, smooth solutions: \emph{FNO}.
\item Sphere / earth-system: \emph{SFNO}.
\item Branch--trunk separation, point queries, branched inputs:
\emph{DeepONet}.
\item Unstructured mesh / complex geometry: \emph{GNO / MGN /
Transolver}.
\item Few-shot transfer across PDE families: \emph{Poseidon, DPOT,
MPP} (foundation operators).
\end{itemize}
\end{insight}
